\documentclass[12pt,a4paper]{article}

\usepackage[spanish,es-nodecimaldot]{babel}
\usepackage[T1]{fontenc}
\usepackage[utf8]{inputenc}
\usepackage{lmodern}
\usepackage{geometry}
\usepackage{setspace}
\usepackage{booktabs}
\usepackage{array}
\usepackage{longtable}
\usepackage{graphicx}
\usepackage{float}
\usepackage{hyperref}
\usepackage{caption}
\usepackage{amsmath}

\geometry{margin=2.5cm}
\onehalfspacing
\hypersetup{colorlinks=true, linkcolor=black, urlcolor=blue, citecolor=black}

\title{\textbf{Primer Proyecto de Inteligencia Artificial 2026}\\
Detección de Ciberacoso en Redes Sociales\\[0.3cm]
\large Opción B --- Clasificación de Tweets}
\author{Ángel Esteban Esquit Hernández -- 23221\\
Javier Eduardo España Pacheco -- 23361\\
Roberto José Barreda Siekavizza -- 23354}
\date{Abril 2026}

\begin{document}

\maketitle
\tableofcontents
\newpage

%% ====================================================================
\section{Resumen Ejecutivo}
%% ====================================================================

Este informe presenta el desarrollo de un sistema de clasificación supervisada
para detectar ciberacoso en tweets. El problema se formula como clasificación
multiclase sobre seis categorías: \textit{age}, \textit{ethnicity},
\textit{gender}, \textit{religion}, \textit{other\_cyberbullying} y
\textit{not\_cyberbullying}. El flujo completo incluye preprocesamiento NLP,
vectorización de texto (BoW, TF-IDF, TF-IDF híbrido y TF-IDF híbrido con
keywords dirigidas), entrenamiento y comparación de tres modelos (Naive Bayes,
Linear SVM y Random Forest), optimización de hiperparámetros mediante
validación cruzada estratificada, análisis de overfitting y pruebas de
usuario con una función interactiva de predicción.

El mejor modelo obtenido es \textbf{Linear SVM (tuned)} con un accuracy de
81.82\% y un F1 ponderado de 0.8182 en el conjunto de prueba.

%% ====================================================================
\section{Presentación y Análisis de Datos}
%% ====================================================================

\subsection{Fuente y estructura del dataset}
Se utilizó el dataset \textit{Cyberbullying Classification} disponible en
Kaggle\footnote{\url{https://www.kaggle.com/datasets/andrewmvd/cyberbullying-classification}},
almacenado localmente en \texttt{data/raw/cyberbullying\_tweets.csv}.
El archivo contiene 47,692 muestras con dos columnas:
\texttt{tweet\_text} (texto del tweet) y \texttt{cyberbullying\_type}
(etiqueta de clasificación).

\subsection{Análisis exploratorio y distribución de clases}
Se verificó la distribución de las seis clases, encontrando un balance
prácticamente uniforme. La clase más numerosa es \textit{religion} con 7,998
muestras y la menor es \textit{other\_cyberbullying} con 7,823; la diferencia
es de apenas 175 muestras (0.37\%), lo que permite enfocarse en el
procesamiento del texto sin necesidad de técnicas de re-muestreo.

\begin{table}[H]
\centering
\caption{Distribución de clases en el dataset (47,692 muestras totales)}
\begin{tabular}{lrr}
\toprule
Clase & Cantidad & Porcentaje \\
\midrule
age                  & 7,992 & 16.76\% \\
ethnicity            & 7,961 & 16.69\% \\
gender               & 7,973 & 16.72\% \\
not\_cyberbullying   & 7,945 & 16.66\% \\
other\_cyberbullying & 7,823 & 16.40\% \\
religion             & 7,998 & 16.77\% \\
\bottomrule
\end{tabular}
\end{table}

\begin{figure}[H]
    \centering
    \includegraphics[width=0.85\textwidth]{data_class_distribution.png}
    \caption{Distribución de clases normalizada (min-max) tras limpieza.}
\end{figure}

\subsection{Datos faltantes y atípicos}
Se realizó una inspección de calidad de datos. No se encontraron valores
faltantes en ninguna de las dos columnas (\texttt{tweet\_text} ni
\texttt{cyberbullying\_type}), por lo que las 47,692 filas se mantuvieron
íntegras. El preprocesamiento de texto (Sección~\ref{sec:preprocessing}) se
encarga de manejar datos atípicos a nivel textual, como tweets que contienen
únicamente URLs o menciones y quedan vacíos tras la limpieza.

%% ====================================================================
\section{Preprocesamiento de Texto (NLP Pipeline)}
\label{sec:preprocessing}
%% ====================================================================

El pipeline de preprocesamiento está implementado en
\texttt{src/text\_preprocessing.py} mediante la clase
\texttt{TextPreprocessor}. Las etapas se aplican secuencialmente:

\subsection{Limpieza del texto}
\begin{itemize}
  \item \textbf{Eliminación de URLs:} se remueven direcciones web
        (http, https, www) mediante expresiones regulares.
  \item \textbf{Eliminación de menciones} (\texttt{@usuario}):
        se eliminan por no correlacionarse con el tipo de ciberacoso.
  \item \textbf{Eliminación de hashtags} (\texttt{\#}):
        se eliminan por completo; su formato concatenado dificulta
        la tokenización.
  \item \textbf{Eliminación de emojis y caracteres especiales:}
        se aplican patrones Unicode para emojis y luego se remueven
        todos los caracteres no alfanuméricos.
\end{itemize}

\subsection{Normalización}
Todo el texto se convierte a minúsculas y se colapsan espacios múltiples
en un solo espacio.

\subsection{Tokenización}
Se utiliza \texttt{word\_tokenize} de NLTK con la opción
\texttt{preserve\_line=True} para dividir cada tweet limpio en tokens.

\subsection{Eliminación de stopwords}
Se eliminan las stopwords del inglés (corpus de NLTK), con la excepción
importante de conservar palabras de negación y pronombres personales:
\texttt{no}, \texttt{nor}, \texttt{not}, \texttt{never}, \texttt{you},
\texttt{your}, \texttt{yours}, \texttt{yourself}, \texttt{yourselves}.
Esta decisión es fundamental porque frases como ``\textit{no one wants you}''
o ``\textit{you are not welcome}'' perderían su significado de acoso si se
eliminaran estas palabras.

\subsection{Lematización}
Se aplica el lematizador WordNet de NLTK para reducir cada token a su forma
base (lema). Por ejemplo, ``running'' $\to$ ``run'' e
``insulted'' $\to$ ``insult''.

\subsection{Ejemplo de transformación completa}
\begin{quote}
\textbf{Original:} \textit{``Go back to your country, we do not want you
here!!! @user123 \#hate''}\\
\textbf{Limpio:} \textit{``go back your country not want you''}
\end{quote}

%% ====================================================================
\section{Vectorización: De Texto a Números}
\label{sec:vectorization}
%% ====================================================================

Los modelos de aprendizaje automático requieren representaciones numéricas.
Se implementaron cuatro esquemas de vectorización en
\texttt{src/vectorization.py} y se compararon objetivamente con el mismo
clasificador de referencia (LinearSVC, $C=0.5$, \texttt{class\_weight=balanced}).

\subsection{Bag of Words (BoW)}
Representa cada documento como un vector de conteos de frecuencia de
unigramas y bigramas (\texttt{ngram\_range=(1,2)}), limitado a 2,000
características (\texttt{max\_features=2000}), con filtros
\texttt{min\_df=2} y \texttt{max\_df=0.95}.

\subsection{TF-IDF}
\textit{Term Frequency -- Inverse Document Frequency} pondera cada término
por su frecuencia local y su rareza global. Se emplea
\texttt{sublinear\_tf=True} para escala logarítmica ($1 + \log(\text{tf})$).
Mismos parámetros de filtrado que BoW.

\subsection{TF-IDF Híbrido (Palabras + Caracteres)}
Combina dos vectorizadores TF-IDF concatenados horizontalmente
(\texttt{hstack}): n-gramas de palabras (1,2) y n-gramas de caracteres a
nivel de palabra (\texttt{char\_wb}, rango 3--5, hasta 1,000 features).
Los n-gramas de caracteres capturan patrones sub-léxicos como sufijos
despectivos, variaciones ortográficas y jerga abreviada.

\subsection{TF-IDF Híbrido + Keywords Dirigidas}
Extiende el esquema híbrido con una tercera capa de características binarias
basadas en un vocabulario fijo de 17 frases de rechazo y xenofobia
(e.g.\ ``go back your country'', ``not welcome here'', ``nobody want you'').
Estas características se multiplican por un peso configurable
(\texttt{keyword\_feature\_weight=3.0}), amplificando la señal de patrones
indicadores de clases difíciles como \textit{other\_cyberbullying}.

\subsection{Comparación de vectorizadores}

\begin{table}[H]
\centering
\caption{Comparación de esquemas de vectorización (mismo clasificador LinearSVC)}
\label{tab:vec_comp}
\begin{tabular}{lccccc}
\toprule
Vectorizador & Accuracy & F1 & Macro F1 & Weak F1 & Tiempo (s) \\
\midrule
tfidf\_hybrid          & 0.8184 & 0.8185 & 0.8175 & 0.5703 & 5.85 \\
tfidf\_hybrid\_keywords & 0.8182 & 0.8182 & 0.8173 & 0.5693 & 6.85 \\
tfidf                  & 0.8108 & 0.8114 & 0.8105 & 0.5630 & 1.44 \\
bow                    & 0.8069 & 0.8087 & 0.8077 & 0.5577 & 1.78 \\
\bottomrule
\end{tabular}
\end{table}

Los esquemas híbridos superan consistentemente a TF-IDF simple y BoW.
Se seleccionó \texttt{tfidf\_hybrid\_keywords} como vectorizador de
producción por su mayor robustez en escenarios de inferencia de usuario,
donde las keywords dirigidas mejoran la clasificación de tweets con
patrones explícitos de rechazo.

%% ====================================================================
\section{Partición de Datos y Metodología}
%% ====================================================================

\subsection{Partición train/test}
Se utilizó \texttt{train\_test\_split} de scikit-learn con partición
80/20 estratificada (\texttt{stratify=y}, \texttt{random\_state=42}):
\begin{itemize}
  \item \textbf{Entrenamiento:} 38,153 muestras (80\%).
  \item \textbf{Prueba:} 9,539 muestras (20\%).
\end{itemize}

\subsection{Validación cruzada para tuning}
Durante la búsqueda de hiperparámetros de Naive Bayes y Linear SVM se
empleó \texttt{StratifiedKFold} con $k=2$ folds sobre el conjunto de
entrenamiento. Se optó por 2 folds para reducir el costo computacional,
dado que cada fold contiene aproximadamente 19,000 muestras.

\subsection{Subconjunto estratificado para Random Forest}
Para controlar el costo computacional del modelo de ensamble, Random Forest
se entrenó sobre un subconjunto estratificado de 12,000 muestras del
conjunto de entrenamiento, reduciendo el tiempo de entrenamiento en
aproximadamente un 68\%.

%% ====================================================================
\section{Modelos Probados y Configuración Técnica}
%% ====================================================================

Se evaluaron tres arquitecturas, cada una representando una familia
distinta de algoritmos.

\subsection{Naive Bayes Multinomial (Modelo Probabilístico)}
Clasificador probabilístico basado en el Teorema de Bayes con suposición
de independencia condicional entre características. Eficiente para
clasificación de texto al operar directamente sobre pesos TF-IDF.

\textbf{Hiperparámetros ajustados:}
\begin{itemize}
  \item \texttt{alpha} (suavizado de Laplace): se probaron
        $\{0.005, 0.01, 0.05\}$. Mejor: \texttt{alpha=0.005}.
  \item Mejor F1 en validación cruzada: 0.7543.
\end{itemize}

\subsection{Linear SVM --- Support Vector Machine (Modelo Lineal)}
\texttt{LinearSVC} implementa una SVM lineal optimizada para alta
dimensionalidad. Busca el hiperplano que maximiza el margen de
separación entre clases.

\textbf{Hiperparámetros ajustados:}
\begin{itemize}
  \item \texttt{C} (regularización): se probaron $\{0.5, 1.0\}$.
        Mejor: $C=0.5$.
  \item \texttt{class\_weight}: se comparó \texttt{None} vs.\
        \texttt{balanced}. Mejor: \texttt{balanced}.
  \item \texttt{max\_iter}: 5,000 (garantizar convergencia).
  \item Mejor F1 en validación cruzada: 0.8109.
\end{itemize}

\subsection{Random Forest (Modelo de Ensamble)}
Construye 120 árboles de decisión sobre submuestras aleatorias
(bagging). Cada árbol se entrena con un subconjunto aleatorio de
características por nodo.

\textbf{Hiperparámetros fijos} (sin tuning por costo computacional):
\begin{itemize}
  \item \texttt{n\_estimators}: 120.
  \item \texttt{max\_depth}: 20.
  \item \texttt{min\_samples\_split}: 2.
  \item Subconjunto de entrenamiento: 12,000 muestras estratificadas.
\end{itemize}

\subsection{Resumen de hiperparámetros}

\begin{table}[H]
\centering
\caption{Hiperparámetros seleccionados y método de tuning}
\label{tab:hyperparams}
\begin{tabular}{p{3.0cm} p{5.5cm} c p{2.5cm}}
\toprule
Modelo & Hiperparámetros & F1-CV & Método \\
\midrule
Naive Bayes   & \texttt{alpha = 0.005}                     & 0.7543 & Grid + CV-2 \\
Linear SVM    & \texttt{C=0.5, class\_weight=balanced}      & 0.8109 & Grid + CV-2 \\
Random Forest & \texttt{n\_est=120, depth=20, split=2}      & N/A    & Fijo \\
\bottomrule
\end{tabular}
\end{table}

%% ====================================================================
\section{Resultados y Comparación de Modelos}
%% ====================================================================

\subsection{Métricas en conjunto de prueba}

\begin{table}[H]
\centering
\caption{Comparación de modelos en test}
\label{tab:model_comp}
\begin{tabular}{lcccccc}
\toprule
Modelo & Accuracy & Precision & Recall & F1 & Weak F1 & Balance \\
\midrule
Linear SVM (tuned)    & 0.8182 & 0.8187 & 0.8182 & 0.8182 & 0.5693 & 0.6938 \\
Random Forest         & 0.7981 & 0.8189 & 0.7981 & 0.8024 & 0.5379 & 0.6702 \\
Naive Bayes (tuned)   & 0.7606 & 0.7563 & 0.7606 & 0.7557 & 0.5083 & 0.6320 \\
\bottomrule
\end{tabular}
\end{table}

La métrica \textit{balance\_score} combina el rendimiento global con
el de las clases más difíciles:
\[
  \text{balance\_score} = 0.5 \times F1_{\text{ponderado}}
  + 0.5 \times \overline{F1}(\text{not\_cb},\; \text{other\_cb})
\]

\subsection{Métricas por clase del mejor modelo}

\begin{table}[H]
\centering
\caption{Reporte por clase: Linear SVM (tuned)}
\label{tab:class_report}
\begin{tabular}{lcccc}
\toprule
Clase & Precision & Recall & F1-Score & Support \\
\midrule
age                  & 0.9645 & 0.9850 & 0.9746 & 1,598 \\
ethnicity            & 0.9731 & 0.9786 & 0.9759 & 1,592 \\
gender               & 0.8775 & 0.8395 & 0.8581 & 1,595 \\
not\_cyberbullying   & 0.5598 & 0.5362 & 0.5477 & 1,589 \\
other\_cyberbullying & 0.5738 & 0.6089 & 0.5908 & 1,565 \\
religion             & 0.9574 & 0.9556 & 0.9565 & 1,600 \\
\midrule
\textbf{Accuracy}    &        &        & \textbf{0.8182} & 9,539 \\
\textbf{Macro avg}   & 0.8177 & 0.8173 & 0.8173 & 9,539 \\
\textbf{Weighted avg}& 0.8187 & 0.8182 & 0.8182 & 9,539 \\
\bottomrule
\end{tabular}
\end{table}

\subsection{Análisis train vs.\ test (overfitting)}

\begin{table}[H]
\centering
\caption{Comparación Train vs.\ Test por modelo (métricas principales)}
\label{tab:overfitting}
\begin{tabular}{llccc}
\toprule
Modelo & Métrica & Train & Test & Gap \\
\midrule
Linear SVM (tuned) & Accuracy      & 0.8729 & 0.8182 & 0.0547 \\
Linear SVM (tuned) & F1            & 0.8731 & 0.8182 & 0.0549 \\
Linear SVM (tuned) & Balance score & 0.7838 & 0.6938 & 0.0901 \\
\midrule
Naive Bayes (tuned) & Accuracy     & 0.7824 & 0.7606 & 0.0219 \\
Naive Bayes (tuned) & F1           & 0.7785 & 0.7557 & 0.0228 \\
Naive Bayes (tuned) & Balance score& 0.6730 & 0.6320 & 0.0409 \\
\midrule
Random Forest       & Accuracy     & 0.8275 & 0.7981 & 0.0294 \\
Random Forest       & F1           & 0.8321 & 0.8024 & 0.0297 \\
Random Forest       & Balance score& 0.7227 & 0.6702 & 0.0525 \\
\bottomrule
\end{tabular}
\end{table}

\textbf{Naive Bayes} presenta el menor gap (F1 gap = 0.023), indicando
que su simplicidad previene el sobreajuste pero también limita su
rendimiento absoluto (\textit{underfitting}).
\textbf{Random Forest} muestra un gap intermedio (F1 gap = 0.030), y
podría beneficiarse de más datos de entrenamiento.
\textbf{Linear SVM} tiene el mayor gap (F1 gap = 0.055), pero sigue
siendo aceptable ya que se encuentra por debajo del umbral típico de
preocupación (0.10), y sus métricas absolutas en test son las más altas.

\subsection{Figuras de soporte}
\begin{figure}[H]
    \centering
    \includegraphics[width=0.88\textwidth]{model_comparison.png}
    \caption{Comparación de métricas por modelo en el conjunto de prueba.}
\end{figure}

\begin{figure}[H]
    \centering
    \includegraphics[width=0.88\textwidth]{train_vs_test_gap.png}
    \caption{Brecha train-test por métrica y modelo.}
\end{figure}

%% ====================================================================
\section{Matriz de Confusión y Análisis de Errores}
%% ====================================================================

\begin{figure}[H]
    \centering
    \includegraphics[width=0.78\textwidth]{confusion_matrix_best.png}
    \caption{Matriz de confusión del Linear SVM (tuned) en el conjunto de prueba.}
\end{figure}

Las cuatro categorías de ciberacoso específico (\textit{age, ethnicity,
gender, religion}) se clasifican con alta precisión (recall $> 83\%$),
mientras que \textit{not\_cyberbullying} y \textit{other\_cyberbullying}
concentran la mayoría de los errores.

\subsection{Confusiones más comunes}

\begin{table}[H]
\centering
\caption{Top 5 confusiones más frecuentes del mejor modelo}
\label{tab:confusions}
\begin{tabular}{llcr}
\toprule
Clase real & Predicha como & Casos & Tasa de error \\
\midrule
not\_cyberbullying   & other\_cyberbullying & 535 & 33.67\% \\
other\_cyberbullying & not\_cyberbullying   & 499 & 31.88\% \\
gender               & other\_cyberbullying & 129 &  8.09\% \\
gender               & not\_cyberbullying   & 107 &  6.71\% \\
other\_cyberbullying & gender               &  88 &  5.62\% \\
\bottomrule
\end{tabular}
\end{table}

La confusión bidireccional entre \textit{not\_cyberbullying} y
\textit{other\_cyberbullying} (535 + 499 = 1,034 errores) es el desafío
principal, causada por la ambigüedad semántica inherente: el acoso genérico
suele ser implícito, sarcástico o sutil y comparte características léxicas
con tweets inofensivos.

\subsection{Ejemplos correctamente clasificados}
\begin{itemize}
  \item \textit{``you idiot death no time anyone you die tomorrow you idiot
        hatred muslim northerner\ldots''}
        $\to$ \textbf{religion} (odio explícito hacia musulmanes).
  \item \textit{``judge turn heat competition you think store remaining
        team''} $\to$ \textbf{not\_cyberbullying} (lenguaje neutral).
\end{itemize}

\subsection{Ejemplos incorrectamente clasificados}
\begin{itemize}
  \item \textit{``vomit''} --- Real: not\_cyberbullying.
        Predicho: other\_cyberbullying.\\
        Análisis: tweet de una sola palabra con connotación negativa;
        el modelo interpreta el término visceral como hostilidad.
  \item \textit{``fact paid 7 fruit cup coffee school insane''} ---
        Real: not\_cyberbullying. Predicho: age.\\
        Análisis: la mención de ``school'' activa patrones de acoso
        por edad; el modelo no distingue una queja cotidiana.
  \item \textit{``saying youre going put old yeller style triggering
        old yeller feeling''} --- Real: gender.
        Predicho: other\_cyberbullying.\\
        Análisis: contenido de género implícito; el modelo no capta la
        referencia cultural.
\end{itemize}

%% ====================================================================
\section{Justificación del Modelo Seleccionado}
%% ====================================================================

Se selecciona \textbf{Linear SVM (tuned)} como modelo final, basado en
el criterio compuesto \textit{balance\_score} que combina el rendimiento
global con el rendimiento en las clases más difíciles.

\subsection{Superioridad cuantitativa}
Linear SVM supera a los demás en todas las métricas principales
(Tabla~\ref{tab:model_comp}): accuracy 81.82\% ($>$ RF 79.81\%
$>$ NB 76.06\%), F1 ponderado 0.8182 ($>$ 0.8024 $>$ 0.7557) y
balance\_score 0.6938 ($>$ 0.6702 $>$ 0.6320).

\subsection{Ventajas adicionales}
\begin{itemize}
  \item \textbf{Eficiencia computacional:} LinearSVC es extremadamente
        rápido. El pipeline completo se ejecutó en 40.15 segundos
        (semilla 42, 1 núcleo).
  \item \textbf{Escalabilidad:} opera eficientemente en el espacio
        disperso de alta dimensionalidad generado por TF-IDF, sin
        requerir reducción de dimensionalidad.
  \item \textbf{Regularización:} $C=0.5$ con \texttt{class\_weight=balanced}
        proporciona un buen equilibrio entre sesgo y varianza.
  \item \textbf{Interpretabilidad:} modelo lineal que facilita la
        explicabilidad, relevante en aplicaciones de moderación de
        contenido.
\end{itemize}

\subsection{Análisis de las clases difíciles}
Las categorías \textit{not\_cyberbullying} y \textit{other\_cyberbullying}
son las más difíciles para todos los modelos. La rama de keywords dirigidas
en la vectorización y la capa de reglas post-modelo en la función de
usuario abordan parcialmente este desafío para escenarios de inferencia.

%% ====================================================================
\section{Pruebas de Usuario}
%% ====================================================================

El sistema incluye una función \texttt{predict\_tweet} en
\texttt{run\_project.py} que recibe un texto libre, aplica preprocesamiento,
vectoriza y retorna la predicción. Adicionalmente, se implementa una
capa transparente de reglas post-modelo para mejorar la clasificación
en escenarios de inferencia directa, sin afectar las métricas del
benchmark.

\begin{table}[H]
\centering
\caption{Resultados de la función de prueba de usuario}
\label{tab:user_tests}
\begin{tabular}{p{5.8cm} p{3.5cm} p{2.8cm} c}
\toprule
Tweet de entrada & Texto limpio & Predicción & Regla \\
\midrule
I hope you have a great day!
  & hope you great day
  & not\_cyberbullying & No \\
\addlinespace
You are stupid and nobody wants you
  & you stupid nobody want you
  & other\_cyberbullying & Sí \\
\addlinespace
Go back to your country, we do not want you here
  & go back your country not want you
  & other\_cyberbullying & Sí \\
\addlinespace
I disagree with your opinion but respect your right to express it
  & disagree your opinion respect your right express
  & not\_cyberbullying & Sí \\
\bottomrule
\end{tabular}
\end{table}

La capa de reglas opera bajo tres condiciones:
(1)~\textbf{patrones seguros} (e.g.\ ``respect your right'') fuerzan
\textit{not\_cyberbullying};
(2)~\textbf{patrones de alerta} (e.g.\ ``go back your country'') activan
\textit{other\_cyberbullying};
(3)~\textbf{rechazo directo + insulto} (e.g.\ ``not want you'' + ``stupid'')
activan \textit{other\_cyberbullying}.
Cuando una regla anula la predicción del modelo, se indica explícitamente
en la salida.

%% ====================================================================
\section{Conclusiones}
%% ====================================================================

Linear SVM (tuned) fue seleccionado como mejor modelo con 81.82\% de
accuracy y F1 de 0.8182, superando a Random Forest y Naive Bayes en
métricas globales y en las clases ambiguas; el principal desafío es la
confusión bidireccional entre \textit{not\_cyberbullying} y
\textit{other\_cyberbullying}, que concentra la mayoría de errores.

\end{document}